% !TEX root = ../main.tex
\chapter{Future Work}
\label{chap:future_work}

Based on the findings and conclusions established in this work, several natural directions emerge for the continuation and expansion of research in Software Engineering and molecular storage systems:

\begin{enumerate}
    \item \textbf{Noisy Channel Modeling with Indels and Substitutions}:
    Extend the stochastic channel model beyond pure erasures to incorporate realistic rates of point mutations (nucleotide substitutions), deletions, and insertions (indels), which are characteristic of chemical synthesis and third-generation sequencing technologies (such as Oxford Nanopore and Pacific Biosciences). Implementing synchronization and alignment algorithms based on Levenshtein distance or HEDGES codes~\cite{press2020} will allow assessing pipeline robustness under phase noise.

    \item \textbf{Integration of Concatenated Error-Correcting Codes}:
    Couple a two-layered architecture to the simulator, combining an inner block error-correcting code (such as Reed--Solomon with 2 bytes per oligo or convolutional codes) applied across each oligonucleotide's seed and payload with the outer rateless LT code, thereby closely aligning the simulator with the full specification employed by \citeonline{erlich2017}.

    \item \textbf{Investigation of Alternative Distributions and Systematic Codes}:
    Explore alternative probability distributions to Luby's classic RSD, such as degree distributions optimized for Raptor codes and Batched Sparse (BATS) codes, as well as systematic fountain code variants where the first $K$ transmitted packets directly correspond to unencoded source blocks, thereby reducing required overhead in small message regimes.

    \item \textbf{Scalability for Large-Scale Datasets}:
    Scale the experimental infrastructure to evaluate megabyte- or gigabyte-scale payloads (expanding the product $K \times L$ through joint increases in block count $K$ and payload size $L$), incorporating concurrent decoding techniques with thread- and process-level parallelism (\texttt{multiprocessing} or OpenMP) and memory-optimized bipartite graph representations.

    \item \textbf{Dedicated Hardware Acceleration (FPGA)}:
    Design and synthesize an IP core for the peeling decoder algorithm using hardware description languages (VHDL or Verilog), targeting low-cost FPGA devices (such as the Digilent Basys 3 board equipped with a Xilinx Artix-7 FPGA). Parallelizing XOR operations and managing adjacency queues in dedicated hardware has the potential to substantially reduce decoding latency, to be quantified following FPGA implementation and synthesis for real-time storage systems.

    \item \textbf{Autonomous and Portable PRNG for Hardware Synthesis}:
    Implement an autonomous, high-speed pseudorandom number generator (such as an arithmetic LFSR with linear shift registers or PCG32/xoshiro family generators), eliminating dependency on the Python interpreter and enabling direct synthesis of neighborhood index regeneration circuits in reconfigurable logic.

    \item \textbf{Goodness-of-Fit Tests on Degree Distribution under Screening}:
    Conduct large-scale sampling trials ($10^5$ or more droplets) utilizing chi-square ($\chi^2$) goodness-of-fit tests and discrete-adapted Kolmogorov--Smirnov tests across the full frequency histograms of each degree $d \in \{1, \ldots, K\}$, deepening the evaluation of any potential statistical association between sequence physicochemical constraints and degree acceptance probabilities in the RSD.
\end{enumerate}
